\documentclass[12pt]{article}
\usepackage[utf8]{inputenc}
\usepackage[T1]{fontenc}
\usepackage{hyperref}
\title{Takpa Language Resources for K-12 Curriculum Development}
\author{Compiled by Claude Sonnet 5 (claude-sonnet-5) for Andr\'as Kornai}
\date{}
\begin{document}
\maketitle

\section{Introduction}
The code tkk (Takpa) is a retired ISO 639-3 code, not a language of its own. Glottolog and IANA record it as retired as a duplicate and merged into Tawang Monpa (twm). Sources give the effective date as 2011-05-18. Resources should therefore be sought under Tawang Monpa (twm) or Dakpa (dka), not under tkk.

\section{Primary Educational Resources}
None applicable to the retired code.

\section{Linguistic and Reference Materials}
\begin{itemize}
\item Glottolog bookkeeping entry: \url{https://glottolog.org/resource/languoid/iso/tkk}
\item IANA subtag record (deprecated, preferred value twm): \url{https://www.iana.org/assignments/lang-subtags-templates/tkk.txt}
\item Wikipedia, Takpa language: \url{https://en.wikipedia.org/wiki/Takpa_language}
\end{itemize}

\section{Religious and Community Resources}
None found.

\section{Digital and Audio Resources}
None found.

\section{Curriculum Development Context}
No curriculum can be built on tkk. Any work should be redirected to Tawang Monpa (twm) or Dakpa (dka).
\end{document}